\scanpage{7}
% archive sheet number: 6

inverses de \deleted{cette} \unsure{dernière|dernier}, par $p_1^0$ \unsure{ou|et} $p_2^0$, sont identiques, car ce sont les obstructions à relever une $=$ section $g_0\in H^0(R_0,\mathcal O_{R_0})$ en un $g\in H^0(R,\mathcal O_R)$. D'où

\[
\begin{aligned}
\unsure{\zeta|\xi}&\in\operatorname{Ker}\bigl(H^1(X_0,\unsure{I|J|\mathcal J}_{X_0})\rightrightarrows H^1(R_0,I_{R_0})\bigr)\\
&=\check H^0(X_0/Y_0,{\unsure{\mathscr H|\mathcal H}}^1(\underline I_{Y_0}))
\end{aligned}
\]

or on sait que le $2^{\mathrm{ième}}$ mb est nul, $X_0\longrightarrow Y_0$ étant couvrant qu. cpct et qu. séparé. Par suite on trouve $\zeta=0$, d'où l'existence d'un $f$. Si on avait \deleted{$f\in\operatorname{Ker}$} $p_1(f)=p_2(f)$,

on aurait gagné, car $f$ proviendrait d'un $\varphi\in A$, qui \deleted{\unsure{redonnerait} \unsure{par}} \inserted{relèverait} donc $\varphi_0$.

\diagram{\begin{tikzcd}[ampersand replacement=\&,column sep=small,row sep=small]
\& f \& {}\\
A\arrow[r]\arrow[d] \& H^0(X,\underline{\mathcal O}_X)\arrow[r,shift left=.6ex,"p_1"]\arrow[r,shift right=.6ex,"p_2"']\arrow[d] \& H^0(R,\mathcal O_R)\arrow[d]\\
A_0\arrow[r] \& H^0(X_0,\mathcal O_{X_0})\arrow[r,shift left=.6ex]\arrow[r,shift right=.6ex] \& H^0(R_0,\mathcal O_{R_0})\\
\varphi_0\& f_0\&{}
\end{tikzcd}}

À priori, on peut seulement dire que $y\unsure{\in|=}p_1(f)-p_2(f)\in H^0(R_0,I_{R_0})$ d'ailleurs on vérifie que

\[
y\in{\unsure{\mathbb Z|Z}}^{1}(X_0/Y_0,\underline I_{Y_0})\qquad\text{i.e. }dy=0
\]

D'ailleurs, on dispose de l'arbitraire de $\lambda\in H^0(X_0,\underline I_{X_0})$ pour $f$, et on voit que \deleted{la correction} dire que $p_1(f+\lambda)=p_{\unsure{2|\,}}(f+\lambda)$ signifie que $p_2(\lambda)-p_1(\lambda)=\deleted{\unsure{u|\text{«}}}y$, i.e.

\[
d(\lambda)=y.
\]
